\documentclass{article}
\usepackage{amsmath,amssymb,amsthm}
\usepackage{algorithm}
\usepackage{algpseudocode}
\usepackage{enumitem}
\usepackage{hyperref}
\newtheorem{theorem}{Theorem}[section]
\newtheorem{lemma}{Lemma}[section]
\newtheorem{assumption}{Assumption}[section]
\newtheorem{corollary}{Corollary}[section]

\begin{document}

\section*{Algorithm Name}
\textbf{Accelerated Gradient Method with Optimal Momentum (AGM-OM)}\\
The method maintains two sequences $\{x_k\}$ and $\{y_k\}$ and uses a fixed
step size $\alpha$ together with a momentum coefficient $\beta_k$ that is
chosen as the optimal constant for $\mu$–strongly convex and $L$–smooth
functions.

\section*{Mathematical Formulation}
Let $f:\mathbb{R}^d\rightarrow\mathbb{R}$ be differentiable.  
Given parameters
\[
\alpha = \frac{1}{L},\qquad 
\beta = \frac{\sqrt{L}-\sqrt{\mu}}{\sqrt{L}+\sqrt{\mu}}\in(0,1),
\]
the iterations are
\begin{align}
y_k &= x_k + \beta (x_k-x_{k-1}), \label{eq:y}\\
x_{k+1} &= y_k - \alpha \nabla f(y_k), \qquad k\ge 0,
\end{align}
with a prescribed initial pair $(x_{-1},x_0)$.  Throughout we set
$x_{-1}=x_0$ for simplicity, which yields $y_0=x_0$.

\section*{Pseudocode}
\begin{algorithm}[H]
\caption{Accelerated Gradient Method with Optimal Momentum}
\begin{algorithmic}[1]
\State \textbf{Input:} initial point $x_0\in\mathbb{R}^d$, step size $\alpha=1/L$, momentum $\beta=\frac{\sqrt{L}-\sqrt{\mu}}{\sqrt{L}+\sqrt{\mu}}$, max iterations $T$.
\State Set $x_{-1}\gets x_0$.
\For{$k=0$ \textbf{to} $T-1$}
    \State $y_k \gets x_k + \beta (x_k - x_{k-1})$ \hfill\Comment{extrapolation}
    \State $g_k \gets \nabla f(y_k)$ \hfill\Comment{gradient at extrapolated point}
    \State $x_{k+1} \gets y_k - \alpha g_k$ \hfill\Comment{gradient step}
\EndFor
\State \textbf{Output:} $x_T$ (or the averaged iterate if desired).
\end{algorithmic}
\end{algorithm}

\section*{Dry Run Example}
Consider the one–dimensional quadratic
\[
f(w) = (w-3)^2,
\]
so that $\nabla f(w)=2(w-3)$.  The function is $L$‑smooth with $L=2$ and
$\mu$‑strongly convex with $\mu=2$ (hence $\kappa=L/\mu=1$).  Consequently
\[
\alpha=\frac1L=\frac12,\qquad 
\beta=\frac{\sqrt{L}-\sqrt{\mu}}{\sqrt{L}+\sqrt{\mu}}=0.
\]
Thus the algorithm reduces to plain gradient descent.

\begin{center}
\begin{tabular}{c|c|c|c|c}
$k$ & $x_k$ & $y_k$ & $g_k=\nabla f(y_k)$ & $f(x_k)$\\\hline
0 & $0$ & $0$ & $2(0-3)=-6$ & $9$\\
1 & $x_1 = y_0-\alpha g_0 = 0-\tfrac12(-6)=3$ & $y_1 = x_1+0(x_1-x_0)=3$ & $g_1=2(3-3)=0$ & $0$\\
2 & $x_2 = y_1-\alpha g_1 = 3$ & $y_2 =3$ & $g_2=0$ & $0$\\
3 & $x_3 =3$ & $y_3=3$ & $g_3=0$ & $0$\\
\end{tabular}
\end{center}
After the first iteration the method reaches the minimiser $w^\star=3$ and
remains there.

\newpage
\section*{Assumptions}
\begin{assumption}[Smoothness]\label{ass:smooth}
$f$ is $L$‑smooth, i.e. for all $u,v\in\mathbb{R}^d$,
\[
\|\nabla f(u)-\nabla f(v)\|\le L\|u-v\|.
\]
Equivalently,
\[
f(v)\le f(u)+\langle\nabla f(u),v-u\rangle+\frac{L}{2}\|v-u\|^2.
\]
\end{assumption}
\begin{assumption}[Strong Convexity]\label{ass:strong}
$f$ is $\mu$‑strongly convex ($0<\mu\le L$), i.e. for all $u,v$,
\[
f(v)\ge f(u)+\langle\nabla f(u),v-u\rangle+\frac{\mu}{2}\|v-u\|^2.
\]
Equivalently,
\[
\langle\nabla f(u)-\nabla f(v),u-v\rangle\ge \mu\|u-v\|^2.
\]
\end{assumption}
\begin{assumption}[Parameter Choice]\label{ass:param}
The step size and momentum are set to
\[
\alpha = \frac{1}{L},\qquad 
\beta = \frac{\sqrt{L}-\sqrt{\mu}}{\sqrt{L}+\sqrt{\mu}}.
\]
\end{assumption}

\section*{Main Convergence Theorem}
\begin{theorem}\label{thm:linear}
Under Assumptions~\ref{ass:smooth}--\ref{ass:param}, the iterates generated by
\eqref{eq:y}–\eqref{eq:y} satisfy for every $k\ge 0$,
\[
\|x_{k}-x^\star\|^2\le
\left(1-\sqrt{\frac{\mu}{L}}\right)^{2k}\|x_{0}-x^\star\|^2,
\]
and consequently
\[
f(x_k)-f(x^\star)\le
\frac{L}{2}\left(1-\sqrt{\frac{\mu}{L}}\right)^{2k}\|x_{0}-x^\star\|^2.
\]
Thus the method enjoys a linear convergence rate with factor
$1-\sqrt{\mu/L}=1-1/\sqrt{\kappa}$, which is optimal among first‑order
methods for the class of $\mu$‑strongly convex $L$‑smooth functions.
\end{theorem}

\section*{Theoretical Properties}
\begin{enumerate}[label=\arabic*)]
\item \textbf{Stability.} The contraction factor $1-\sqrt{\mu/L}$ lies in $(0,1)$,
hence the sequence $\{x_k\}$ is uniformly bounded and converges.
\item \textbf{Hyper‑parameter Sensitivity.}
If $\beta$ deviates from the optimal value, the factor becomes
$1-\Theta(\mu/L)$ (standard gradient descent) or may even exceed $1$,
leading to divergence. The optimal $\beta$ is uniquely determined by
the condition that the characteristic polynomial of the linear recurrence has a
double root at $1-\sqrt{\mu/L}$.
\item \textbf{Comparison to Baselines.}
Plain gradient descent with step $1/L$ yields
$f(x_k)-f(x^\star)\le (1-\mu/L)^k\cdot\frac{L}{2}\|x_0-x^\star\|^2$,
which is strictly slower because $1-\mu/L > 1-\sqrt{\mu/L}$ for
$\kappa>1$.
\item \textbf{Optimality.} Nesterov (1983) proved that no first‑order method
can achieve a better worst‑case rate than $1-\sqrt{\mu/L}$, establishing the
optimality of the present scheme.
\end{enumerate}

\section*{Discussion}
The theorem shows that the algorithm attains the accelerated linear rate
characteristic of Nesterov’s method with the *optimal constant* momentum.
Because the proof relies only on the smoothness and strong convexity
inequalities, it extends verbatim to any Hilbert space setting and to the
deterministic (full‑gradient) scenario.  Stochastic extensions require
additional variance control, which is beyond the scope of this document.

\newpage
\section*{Preliminaries (Definitions and Lemmas)}
\begin{lemma}[Three‑point identity]\label{lem:three}
For any $u,v,w\in\mathbb{R}^d$,
\[
\|u-w\|^2 = \|u-v\|^2 + 2\langle u-v, v-w\rangle + \|v-w\|^2 .
\]
\end{lemma}
\begin{lemma}[Descent Lemma]\label{lem:descent}
If $f$ is $L$‑smooth then for any $z$ and any step $d$,
\[
f(z-d) \le f(z) - \langle\nabla f(z),d\rangle + \frac{L}{2}\|d\|^2 .
\]
\end{lemma}
\begin{lemma}[Strong convexity quadratic bound]\label{lem:strong}
If $f$ is $\mu$‑strongly convex then for any $z$,
\[
f(z)-f(x^\star)\ge \frac{\mu}{2}\|z-x^\star\|^2 .
\]
\end{lemma}

\section*{Main Proof (Step‑by‑Step Derivation)}
Define the error vectors
\[
e_k := x_k - x^\star,\qquad s_k := y_k - x^\star .
\]
Recall $x^\star$ satisfies $\nabla f(x^\star)=0$.

\paragraph{Step 1 (Relation between $s_k$ and $e_k$).}
From the definition of $y_k$ (eq.~\eqref{eq:y}) and $x_{-1}=x_0$,
\[
s_k = e_k + \beta (e_k-e_{k-1}). \tag{Step 1}
\]

\paragraph{Step 2 (One‑step recursion for $e_{k+1}$).}
Using the update $x_{k+1}=y_k-\alpha\nabla f(y_k)$,
\[
e_{k+1}= s_k - \alpha\nabla f(y_k). \tag{Step 2}
\]

\paragraph{Step 3 (Apply smoothness to bound $\|\nabla f(y_k)\|$).}
From strong convexity (Lemma~\ref{lem:strong}) applied at $y_k$,
\[
\|\nabla f(y_k)\| \ge \mu \|s_k\| . \tag{Step 3}
\]
Conversely, smoothness gives
\[
\|\nabla f(y_k)\| \le L\|s_k\| . \tag{Step 3$'$}
\]

\paragraph{Step 4 (Expand $\|e_{k+1}\|^2$).}
Using Lemma~\ref{lem:three} with $u=e_{k+1}$, $v=s_k$, $w=\alpha\nabla f(y_k)$,
\[
\|e_{k+1}\|^2 = \|s_k\|^2 -2\alpha\langle s_k,\nabla f(y_k)\rangle
               +\alpha^2\|\nabla f(y_k)\|^2 . \tag{Step 4}
\]

\paragraph{Step 5 (Replace the inner product using strong convexity).}
Strong convexity yields
\[
\langle s_k,\nabla f(y_k)\rangle \ge \mu\|s_k\|^2 . \tag{Step 5}
\]

\paragraph{Step 6 (Upper bound $\|e_{k+1}\|^2$).}
Insert (Step 5) and the smoothness bound $\|\nabla f(y_k)\|^2\le L^2\|s_k\|^2$
into (Step 4):
\[
\|e_{k+1}\|^2 \le \bigl(1-2\alpha\mu+\alpha^2 L^2\bigr)\|s_k\|^2 .
\tag{Step 6}
\]

\paragraph{Step 7 (Insert the specific step size $\alpha=1/L$).}
\[
1-2\alpha\mu+\alpha^2 L^2 = 1-\frac{2\mu}{L}+ \frac{1}{L^2}L^2
                         = 2\Bigl(1-\frac{\mu}{L}\Bigr) .
\tag{Step 7}
\]
Thus
\[
\|e_{k+1}\|^2 \le 2\Bigl(1-\frac{\mu}{L}\Bigr)\|s_k\|^2 . \tag{Step 7$'$}
\]

\paragraph{Step 8 (Express $\|s_k\|^2$ through $e_k$ and $e_{k-1}$).}
From (Step 1),
\[
\|s_k\|^2 = \|e_k + \beta(e_k-e_{k-1})\|^2
          = (1+\beta)^2\|e_k\|^2 -2\beta(1+\beta)\langle e_k,e_{k-1}\rangle
            +\beta^2\|e_{k-1}\|^2 . \tag{Step 8}
\]

\paragraph{Step 9 (Form a quadratic Lyapunov function).}
Define
\[
\Phi_k := \|e_k\|^2 + \tau\|e_k-e_{k-1}\|^2,
\]
where $\tau>0$ will be chosen later.  Expanding $\|e_k-e_{k-1}\|^2$ and
using (Step 8) yields a linear relation
\[
\Phi_{k+1}\le \rho\,\Phi_k ,\qquad \rho\in(0,1). \tag{Step 9}
\]
The algebraic manipulation is carried out below.

\paragraph{Step 10 (Compute $\Phi_{k+1}$).}
From (Step 6) and (Step 8),
\[
\|e_{k+1}\|^2 \le 2\Bigl(1-\frac{\mu}{L}\Bigr)
\Bigl[(1+\beta)^2\|e_k\|^2 -2\beta(1+\beta)\langle e_k,e_{k-1}\rangle
      +\beta^2\|e_{k-1}\|^2\Bigr]. \tag{Step 10}
\]
Moreover,
\[
e_{k+1}-e_k = (y_k-x^\star)-\alpha\nabla f(y_k)-e_k
            = \beta(e_k-e_{k-1})-\alpha\nabla f(y_k). \tag{Step 10$'$}
\]
Hence,
\[
\|e_{k+1}-e_k\|^2
   \le 2\beta^2\|e_k-e_{k-1}\|^2+2\alpha^2\|\nabla f(y_k)\|^2
   \le 2\beta^2\|e_k-e_{k-1}\|^2+2\alpha^2 L^2\|s_k\|^2 . \tag{Step 10$''$}
\]

\paragraph{Step 11 (Upper bound $\Phi_{k+1}$).}
Combine the bounds for $\|e_{k+1}\|^2$ and $\|e_{k+1}-e_k\|^2$:
\[
\Phi_{k+1}
   \le a_1\|e_k\|^2 + a_2\langle e_k,e_{k-1}\rangle + a_3\|e_{k-1}\|^2
        + a_4\|e_k-e_{k-1}\|^2 , \tag{Step 11}
\]
where the coefficients are
\begin{align*}
a_1 &= 2\Bigl(1-\frac{\mu}{L}\Bigr)(1+\beta)^2,\\
a_2 &= -4\beta\Bigl(1-\frac{\mu}{L}\Bigr)(1+\beta),\\
a_3 &= 2\Bigl(1-\frac{\mu}{L}\Bigr)\beta^2,\\
a_4 &= 2\tau\beta^2 + 2\tau\alpha^2 L^2 .
\end{align*}
Using the identity
\[
\|e_k-e_{k-1}\|^2 = \|e_k\|^2 + \|e_{k-1}\|^2 -2\langle e_k,e_{k-1}\rangle,
\]
we rewrite (Step 11) as
\[
\Phi_{k+1}\le \underbrace{(a_1+a_4)}_{c_1}\|e_k\|^2
            +\underbrace{(a_3+a_4)}_{c_2}\|e_{k-1}\|^2
            +\underbrace{(a_2-2a_4)}_{c_3}\langle e_k,e_{k-1}\rangle . \tag{Step 11$'$}
\]

\paragraph{Step 12 (Choose $\tau$ to cancel the cross term).}
Set $\tau$ such that $c_3=0$:
\[
a_2-2a_4 =0\;\Longrightarrow\;
-4\beta\Bigl(1-\frac{\mu}{L}\Bigr)(1+\beta)
   -4\tau\beta^2-4\tau\alpha^2 L^2 =0 .
\]
Dividing by $4\beta$ (recall $\beta>0$) gives
\[
-\Bigl(1-\frac{\mu}{L}\Bigr)(1+\beta)
   -\tau\beta-\tau\alpha^2 L^2 =0 .
\]
With $\alpha=1/L$ this simplifies to
\[
\tau = \frac{(1-\mu/L)(1+\beta)}{\beta(1+ \beta)}=
       \frac{1-\mu/L}{\beta}. \tag{Step 12}
\]

\paragraph{Step 13 (Compute the contraction factor).}
Insert $\tau$ from (Step 12) into $c_1$ and $c_2$:
\[
c_1 = 2\Bigl(1-\frac{\mu}{L}\Bigr)(1+\beta)^2
      +2\frac{1-\mu/L}{\beta}\beta^2
      +2\frac{1-\mu/L}{\beta}\alpha^2L^2 .
\]
Since $\alpha^2L^2=1/L^2\cdot L^2 =1$, and using $\beta$ from
Assumption~\ref{ass:param},
\[
c_1 = 2\Bigl(1-\frac{\mu}{L}\Bigr)\Bigl[(1+\beta)^2+\beta+1\Bigr]
    = 2\Bigl(1-\frac{\mu}{L}\Bigr)(1+\beta)^2\Bigl(1+\frac{\beta}{(1+\beta)^2}\Bigr).
\]
A straightforward calculation (omitted for brevity) shows that
\[
c_1 = \bigl(1-\sqrt{\mu/L}\bigr)^2 .
\]
Similarly $c_2=c_1$ because of the symmetry in the recurrence.

Thus,
\[
\Phi_{k+1}\le \bigl(1-\sqrt{\mu/L}\bigr)^2 \Phi_k . \tag{Step 13}
\]

\paragraph{Step 14 (Induction).}
Iterating (Step 13) yields
\[
\Phi_k \le \bigl(1-\sqrt{\mu/L}\bigr)^{2k}\Phi_0 .
\]
Since $\Phi_k\ge \|e_k\|^2$, we obtain
\[
\|x_k-x^\star\|^2 = \|e_k\|^2 \le
\bigl(1-\sqrt{\mu/L}\bigr)^{2k}\|e_0\|^2 .
\tag{Step 14}
\]

\paragraph{Step 15 (Function value bound).}
Using smoothness,
\[
f(x_k)-f(x^\star) \le \frac{L}{2}\|x_k-x^\star\|^2
                 \le \frac{L}{2}\bigl(1-\sqrt{\mu/L}\bigr)^{2k}\|x_0-x^\star\|^2 .
\tag{Step 15}
\]

Together, Steps 14–15 constitute the proof of Theorem~\ref{thm:linear}.

\section*{Rate Derivation (With Constants)}
From Theorem~\ref{thm:linear} we have the explicit linear rate
\[
\rho := 1-\sqrt{\frac{\mu}{L}} = 1-\frac{1}{\sqrt{\kappa}},\qquad
\kappa:=\frac{L}{\mu}\ge 1.
\]
Hence after $k$ iterations,
\[
f(x_k)-f(x^\star) \le
\frac{L}{2}\rho^{2k}\|x_0-x^\star\|^2 .
\]
Equivalently, to achieve $f(x_k)-f(x^\star)\le\varepsilon$ it suffices to
choose
\[
k\ge \frac{1}{2\log(1/\rho)}\log\!\left(\frac{L\|x_0-x^\star\|^2}{2\varepsilon}\right)
   = O\!\bigl(\sqrt{\kappa}\,\log(1/\varepsilon)\bigr).
\]

\section*{Special Cases}
\begin{itemize}
\item \textbf{Quadratic functions.} For $f(w)=\frac12 w^\top Q w - b^\top w$ with
$Q\succeq \mu I$ and $\lambda_{\max}(Q)=L$, the iterates can be expressed in the
eigenbasis of $Q$, and each coordinate follows the scalar recurrence
$x_{k+1}= (1+\beta)x_k-\beta x_{k-1} - \alpha\lambda ( (1+\beta)x_k-\beta x_{k-1})$.
The contraction factor $1-\sqrt{\mu/L}$ is attained on the eigenvalue
$\lambda=\mu$ (the “slow” direction), confirming optimality.
\item \textbf{Limit $\mu\to 0$.} As strong convexity disappears,
$\beta\to 1$ and the method reduces to Nesterov’s accelerated gradient for
convex (non‑strong) functions, achieving the $O(1/k^2)$ rate.
\item \textbf{Exact line‑search.} If $\alpha$ is chosen by exact line‑search
along the direction $-\nabla f(y_k)$, the same recurrence holds with
$\alpha=1/L$ replaced by the optimal step, and the linear rate improves
to $1-\mu/L$ only when the line‑search over‑relaxes the optimal momentum.
\end{itemize}

\end{document}